\documentclass{article}

% if you need to pass options to natbib, use, e.g.:
% \PassOptionsToPackage{numbers, compress}{natbib}
\usepackage{neurips_2025}

\usepackage[utf8]{inputenc}
\usepackage[T1]{fontenc}
\usepackage{hyperref}
\usepackage{url}
\usepackage{booktabs}
\usepackage{amsfonts}
\usepackage{nicefrac}
\usepackage{microtype}
\usepackage{xcolor}
\usepackage{amsmath}
\usepackage{graphicx}

\title{Future-Seed: Cross-Layer State Feedback for In-Place Constraint Repair}

\author{
  Anonymous Authors \\
  Paper under double-blind review
}

\begin{document}
\maketitle

\begin{abstract}
Recurrent and state-space language models are attractive for long-context processing, but they are structurally weak at \emph{in-place constraint repair}: early masked positions must often depend on information located to their right, while one-shot parallel infill easily violates global consistency. We study \textbf{Future-Seed}, a lightweight cross-layer state feedback mechanism that initializes each layer with the previous layer's terminal state, effectively creating a ``re-read across depth'' pathway without changing the causal scan at decode time. On hard synthetic tasks, Future-Seed consistently improves future-aware repair, turns KVSORT from exact failure to exact recovery, and materially improves PERMFILL length extrapolation when paired with a tiny anchor. We then audit the mechanism in a single-GPU post-training campaign across protein secondary-structure spot labeling, Hotpot- and SQuAD-derived text-restoration probes, MBPP long-context code probes, ARC multiple-choice token probes, and punctuation restoration. All real-task numbers are teacher-forced token-accuracy deltas rather than benchmark end metrics. The strongest repeatable real-task family is protein secondary-structure spot labeling, with smaller positive pockets on Hotpot-derived text restoration, punctuation restoration, and SQuAD-derived text restoration. By contrast, MBPP-longctx and ARC token probes show promising but unstable gains and fail strict held-out confirmation. The combined evidence supports Future-Seed as a targeted mechanism for constraint repair rather than a universal post-training booster.
\end{abstract}

\section{Introduction}
Modern sequence models often need to edit an existing sequence rather than generate a new one from scratch. This includes infilling code, restoring structure, repairing permutations, or solving puzzles with hard constraints. These \emph{in-place} tasks are difficult because the masked region must depend on evidence that may only appear on the right side of the input, while the final prediction must remain globally consistent.

This regime overlaps with span corruption and iterative refinement settings that have been studied in text-to-text transfer, non-autoregressive generation, and edit-based decoding~\citep{raffel2020t5,lee2018iterative,ghazvininejad2019maskpredict,gu2019levenshtein}. Our focus is narrower: we ask what happens when the backbone itself remains recurrent and causal at decode time, but deeper layers receive a compact summary of future context through cross-layer state feedback.

Autoregressive Transformers~\citep{vaswani2017attention} and bidirectional masked models such as BERT~\citep{devlin2019bert} provide two different approaches to this problem, but recurrent models remain attractive when prompt length, memory, and streaming efficiency matter. Recent recurrent and state-space model families, including RWKV~\citep{peng2023rwkv} and Mamba~\citep{gu2024mamba}, still inherit a basic weakness: if the model processes tokens left-to-right, early states only see future information through whatever summaries survive in the recurrent state.

This paper studies a minimal intervention, \textbf{Future-Seed}, in which the terminal state of layer $\ell-1$ becomes the initial state of layer $\ell$. The mechanism is tiny, architecture-local, and decode-compatible, but it creates an additional path for future information to influence early masked positions. We ask two practical questions: (1) does this mechanism genuinely help future-aware repair, and (2) does any of that signal survive when we move from toy constraints to real post-training tasks?

\paragraph{Contributions.}
We make three contributions. First, we present a compact mechanism view of Future-Seed as cross-layer state feedback for in-place constraint repair. Second, we show strong synthetic evidence on future-aware sanity tasks, permutation repair, and Sudoku. Third, we report a large single-GPU post-training audit with both positive and negative findings: \texttt{protein\_ss\_spot} is the strongest repeatable real-task family, while \texttt{hotpot\_text\_restore}, \texttt{punc\_restore}, and \texttt{squad\_text\_restore} show smaller positive pockets and \texttt{mbpp\_longctx\_probe}/\texttt{arc\_mc\_probe} remain mixed.

\section{Method}
Let $h_t^{(\ell)}$ denote the hidden state of layer $\ell$ after processing token position $t$. In a standard recurrent stack, each layer starts from its own default initial state. Future-Seed instead injects the previous layer's terminal state into deeper layers:
\begin{equation}
    h_0^{(\ell)} = g^{(\ell)} \odot h_T^{(\ell-1)}, \quad \ell \geq \ell_{\mathrm{start}},
\end{equation}
where $g^{(\ell)}$ is a scalar or head-wise gate and $T$ is the sequence length. Decode remains causal; the mechanism only changes how deeper layers are initialized while scanning the prompt.

\begin{figure}[t]
    \centering
    \fbox{\begin{minipage}{0.93\linewidth}
    \small
    \textbf{Baseline:} each layer starts from its own empty initial state and scans left-to-right once.\\[2mm]
    \textbf{Future-Seed:} layer $\ell$ receives the terminal summary from layer $\ell-1$, so early positions in deeper layers can condition on information already aggregated from the whole prompt.\\[2mm]
    In practice, this acts like a cheap ``re-read across depth'' path for future-aware masked repair.
    \end{minipage}}
    \caption{Mechanism sketch. Future-Seed preserves the causal scan but adds cross-layer state feedback.}
    \label{fig:mechanism}
\end{figure}

This design targets a specific failure mode: masked regions whose correct values depend on suffix information and whose outputs must satisfy a global constraint, such as permutation validity or Sudoku consistency. It is not intended as a universal replacement for attention or iterative refinement.

\section{Experimental Setup}
\paragraph{Backbone and compute.}
All experiments use an RWKV7 backbone with prompt-time Future-Seed only, initialized from the repository checkpoint \texttt{rwkv7-g1d-0.1b-20260129-ctx8192.pth}. Decode remains left-to-right and causal. The real-task campaign is run on a single NVIDIA RTX 4090, using serial quick-to-medium screening to keep wall-clock budget bounded.

\paragraph{Synthetic tasks.}
We evaluate RIGHTCOPY and CONSTR as minimal future-aware probes. We use KVSORT for in-place permutation repair, PERMFILL for anchored permutation filling with length extrapolation, and 4$\times$4 Sudoku for structured constraint repair. Together these tasks span copy, permutation, and grid-style consistency repair.

\paragraph{Real-task audit.}
The real-task campaign covers protein secondary-structure spot labeling, Hotpot- and SQuAD-derived text restoration, long-context MBPP code probes, ARC multiple-choice token probes, punctuation restoration, and a set of lower-ROI constraint proxies. The late-stage queues primarily compare scalar or head-wise prompt gates starting at layers 6, 8, or 10, with common candidates such as \texttt{scalar\_l8\_train1e4}, \texttt{scalar\_l8\_sched\_cos}, \texttt{head\_l8}, and depth-shifted variants.

\paragraph{Queue policy.}
The final fastdiscover queues use a fixed two-stage policy. Quick stage budgets are task-dependent but typically 140--160 steps; medium stage budgets are typically 260--280 steps. We prune candidates below $-0.5$pp at quick stage, promote only the top candidate above $+0.8$pp, and then report medium-stage deltas against the same-task baseline.

\paragraph{Evaluation summary.}
For synthetic tasks we report direct task metrics such as token accuracy, exact match, and solve-rate. For the post-training audit we report teacher-forced answer/token-accuracy deltas in percentage points (pp) by probe family, together with positive-count summaries and qualitative judgments.

\section{Synthetic Mechanism Evidence}
\begin{table}[t]
\centering
\small
\caption{Synthetic mechanism evidence. Future-Seed consistently improves future-aware repair tasks and turns KVSORT from failure to exact recovery.}
\label{tab:synthetic-main}
\begin{tabular}{lccc}
\toprule
Task / metric & Baseline & Future-Seed & Delta \\
\midrule
RIGHTCOPY acc & 10.46\% & 15.50\% & +5.04pp \\
CONSTR acc & 9.77\% & 21.14\% & +11.37pp \\
KVSORT id exact & 0.00\% & 100.00\% & +100.00pp \\
KVSORT ood exact & 0.00\% & 100.00\% & +100.00pp \\
Sudoku solve @ holes=12 & 0.00\% & 55.10\% & +55.10pp \\
\bottomrule
\end{tabular}
\end{table}


Table~\ref{tab:synthetic-main} shows the clearest mechanism evidence. Future-Seed improves both future-aware sanity tasks and produces a sharp failure-to-success transition on KVSORT. The Sudoku result shows that the mechanism is not restricted to one-dimensional copy tasks: it also helps a small structured constraint problem where row, column, and block consistency must be maintained jointly.

\begin{table}[t]
\centering
\small
\caption{PERMFILL length extrapolation. A tiny anchor in the masked span materially improves out-of-distribution exact match.}
\label{tab:permfill-main}
\resizebox{\linewidth}{!}{%
\begin{tabular}{lccccc}
\toprule
Weights & Anchor & n=24 & n=28 & n=32 & n=36 \\
\midrule
FS=1, train n=24 & off & 100.00\% & 0.00\% & 0.00\% & 0.00\% \\
FS=1, train n=24 & k=2 & 99.00\% & 83.00\% & 12.00\% & 0.00\% \\
FS=1, train n=32 & k=2 & 100.00\% & 100.00\% & 100.00\% & 93.50\% \\
\bottomrule
\end{tabular}}
\end{table}


PERMFILL reveals a complementary point. Future-Seed alone is not the whole story for out-of-distribution extrapolation, but a tiny anchor inside the masked span dramatically improves stability. This suggests that the mechanism is especially useful when the problem supplies at least a minimal alignment foothold and the model must propagate global consistency through depth rather than reconstruct structure from a fully unconstrained masked segment.

The synthetic suite also exposes a negative result: simply restoring bidirectional visibility with Transformer baselines does not solve KVSORT under the tested setup. Appendix~A shows that Hungarian decoding can restore validity, but not ordering, which strengthens the case that Future-Seed is doing more than acting as a trivial decoder-side constraint patch.

\section{Real-Task Post-Training Evidence}
\begin{table*}[t]
\centering
\small
\caption{Post-training summary across task families. Counts aggregate medium-stage Future-Seed vs. baseline comparisons from the full search archive.}
\label{tab:posttrain-main}
\begin{tabular}{lcccc}
\toprule
Task family & Best gain & Median gain & Positive med count & Current judgment \\
\midrule
protein\_ss & +8.14pp & +2.18pp & 103/126 & strongest repeatable \\
hotpot & +4.20pp & +0.54pp & 35/53 & repeatable positive \\
mbpp\_longctx & +10.00pp & +0.91pp & 53/87 & promising, unconfirmed \\
arc\_mc & +20.83pp & +4.17pp & 70/116 & high upside, unstable \\
squad & +7.31pp & +0.55pp & 24/35 & mixed, not locked \\
punc & +2.16pp & +0.36pp & 16/24 & small support \\
\bottomrule
\end{tabular}
\end{table*}


Table~\ref{tab:posttrain-main} summarizes the curated post-training archive used by this paper. These counts come from a sequential search archive, not from a balanced i.i.d. trial design, so they should be read as a stability audit rather than a classical multi-seed benchmark. All real-task rows are teacher-forced token-accuracy probes rather than benchmark end metrics such as EM, F1, or pass@k. Two patterns matter.

First, the mechanism is \emph{not} a universal gain switch. The strongest repeatable family is \texttt{protein\_ss\_spot}, which shows both a strong best case and a positive median over a large number of medium-stage comparisons. \texttt{hotpot\_text\_restore} is weaker in absolute magnitude but still consistently positive enough to remain a credible supporting family. \texttt{punc\_restore} and \texttt{squad\_text\_restore} also show small positive pockets.

Second, the most exciting spikes are not the most trustworthy claims. \texttt{mbpp\_longctx\_probe} and \texttt{arc\_mc\_probe} both produce large wins in parts of the search, but neither survives a strict held-out confirmation pass. In other words, Future-Seed can create real-task upside on token-level probes, but the strongest real-task claims are still the ones that survive repeated queue pressure rather than the largest one-off gains.

We therefore frame the real-task evidence conservatively. The paper can defend \texttt{protein\_ss\_spot} as the strongest repeatable family under the current recipe, and \texttt{hotpot\_text\_restore}/\texttt{punc\_restore}/\texttt{squad\_text\_restore} as modest supporting signals. It cannot honestly claim stable held-out confirmation on \texttt{mbpp\_longctx\_probe} or \texttt{arc\_mc\_probe} yet.

\section{Limitations and Negative Results}
This work has four important limitations. First, the post-training campaign is compute-constrained: all decisions are made under a single-GPU search budget, so the resulting archive reflects a pragmatic search policy rather than a clean, balanced hyperparameter study. Second, the aggregated family statistics in Table~\ref{tab:posttrain-main} are not independent repeated trials; they summarize a sequential exploration process. Third, the positive real-task evidence is uneven across families, and the strict confirmation queues fail for the most ambitious real-task headlines. Fourth, the synthetic tasks are intentionally sharp and mechanistic, which helps isolate the effect of Future-Seed but also narrows the external validity of the claim. The anonymous repository snapshot also does not include the referenced RWKV7 checkpoint blob or a pinned training environment, so it is an auditable artifact snapshot rather than a fully turnkey rerun package.

There are still two useful lessons in these negatives. One is scientific: Future-Seed appears better interpreted as a targeted repair mechanism than as a general-purpose prompt-time booster. The other is methodological: aggressive quick-to-medium pruning is valuable because it exposes where a mechanism repeatedly fails before too much budget is spent chasing unstable spikes.

\paragraph{Broader impacts.}
The positive side of this work is better controllable structured editing: code completion, constrained infill, and sequence repair all benefit from mechanisms that preserve global consistency. The risk is that stronger infill and repair mechanisms could also help misuse cases such as stealthier document editing, constraint-satisfying synthetic content, or automated puzzle/benchmark overfitting. We therefore recommend releasing the mechanism together with explicit limitations, exact reproduction commands, and a clear separation between strong synthetic evidence and mixed real-task evidence.

\section{Related Work}
Future-Seed sits between three lines of work. The first is recurrent or state-space sequence modeling, including RWKV~\citep{peng2023rwkv} and Mamba~\citep{gu2024mamba}. The second is masked, span-corruption, and edit-based generation, including BERT~\citep{devlin2019bert}, T5~\citep{raffel2020t5}, iterative refinement~\citep{lee2018iterative}, Mask-Predict~\citep{ghazvininejad2019maskpredict}, Levenshtein Transformer~\citep{gu2019levenshtein}, and MaskGIT~\citep{chang2022maskgit}. The third is benchmark families that stress multi-hop reasoning, code completion, and protein sequence modeling, including HotpotQA~\citep{yang2018hotpotqa}, ARC~\citep{clark2018arc}, SQuAD~\citep{rajpurkar2016squad}, MBPP~\citep{austin2021program}, TAPE~\citep{rao2019tape}, and large-scale protein language modeling~\citep{rives2021biological}. Our contribution is not a new benchmark or a new backbone, but a controlled study of one lightweight cross-layer mechanism across both synthetic repair tasks and a realistic post-training search archive.

\section{Conclusion}
Future-Seed adds a small amount of cross-layer state feedback to a recurrent stack, but its effect is highly structured. The synthetic evidence is strong: future-aware repair, permutation repair, and Sudoku all improve materially. The real-task evidence is narrower but still useful: \texttt{protein\_ss\_spot} is a robust positive family, while \texttt{hotpot\_text\_restore}, \texttt{punc\_restore}, and \texttt{squad\_text\_restore} provide supporting gains and \texttt{mbpp\_longctx\_probe}/\texttt{arc\_mc\_probe} remain mixed. The most defensible conclusion is therefore targeted rather than universal: Future-Seed is a practical mechanism for in-place constraint repair, and a promising but still unstable ingredient for broader post-training.

\bibliographystyle{plainnat}
\bibliography{references}

\clearpage
\appendix
\section{Additional Experimental Details}
\paragraph{Base setting.}
All experiments use an RWKV7 backbone with prompt-time Future-Seed only: generation remains left-to-right and causal at decode time. The real-task campaign uses a single NVIDIA RTX 4090 and serial quick-to-medium screening to conserve wall-clock budget.

\paragraph{Post-training search policy.}
The late-stage queues use a two-stage policy: a quick stage for baseline and candidate screening, and a medium stage only for candidates that clear a positive gate. In the final fastdiscover queues, we prune candidates below $-0.5$pp at quick stage and only promote a top candidate above $+0.8$pp to medium budget.

\paragraph{Task families.}
The real-task audit covers protein secondary-structure labeling, multi-hop QA, code-generation-style long-context completion, ARC multiple choice, punctuation restoration, and a small set of constraint and planning proxies. The most stable family under this setup is protein secondary-structure labeling.

\section{Additional Synthetic Results}
\begin{table*}[t]
\centering
\small
\caption{Transformer baselines on KVSORT. Hungarian decoding restores validity but not ordering.}
\label{tab:appendix-transformer}
\begin{tabular}{lcccc}
\toprule
Model & Decode & Exact & Key-valid & Key-order \\
\midrule
Transformer-MLM & argmax & 0.00\% & 0.00\% & 0.00\% \\
Transformer-MLM & hungarian & 0.00\% & 100.00\% & 0.00\% \\
Transformer-Causal & hungarian & 0.00\% & 100.00\% & 0.00\% \\
Transformer-Causal + ATTN\_FS & hungarian & 0.00\% & 100.00\% & 0.00\% \\
Transformer-MLM + Sinkhorn & hungarian & 0.00\% & 100.00\% & 0.00\% \\
\bottomrule
\end{tabular}
\end{table*}

\begin{table}[t]
\centering
\small
\caption{4x4 Sudoku solve-rate phase curve. Future-Seed delays the collapse point to harder puzzles.}
\label{tab:appendix-sudoku-phase}
\begin{tabular}{rcc}
\toprule
Holes & FS=0 solve & FS=1 solve \\
\midrule
4 & 35.30\% & 100.00\% \\
6 & 16.65\% & 100.00\% \\
8 & 5.20\% & 100.00\% \\
10 & 0.95\% & 95.10\% \\
12 & 0.00\% & 55.10\% \\
14 & 0.00\% & 0.00\% \\
\bottomrule
\end{tabular}
\end{table}

\begin{table}[t]
\centering
\small
\caption{A simple Sudoku consistency regularizer helps at holes=12 but does not solve the hardest setting.}
\label{tab:appendix-sudoku-consistency}
\begin{tabular}{lccc}
\toprule
Setting & holes=10 & holes=12 & holes=14 \\
\midrule
lambda=0.0 & 93.85\% & 47.95\% & 0.00\% \\
lambda=0.1 & 95.15\% & 58.55\% & 0.00\% \\
\bottomrule
\end{tabular}
\end{table}


\paragraph{Takeaway.}
The synthetic suite supports a narrow but strong claim: Future-Seed is especially helpful when a recurrent model must repair a masked region using information that arrives on the right and must also satisfy a global consistency constraint. The Transformer baselines in Table~\ref{tab:appendix-transformer} show that simply restoring bidirectional visibility or decoding with Hungarian matching is not enough on KVSORT; the useful signal comes from the specific cross-layer state feedback path.

\section{Additional Post-Training Results}
\begin{table}[t]
\centering
\small
\caption{Additional task families discussed in appendix.}
\label{tab:appendix-negative-tasks}
\begin{tabular}{lcccc}
\toprule
Task family & Best gain & Median gain & Positive med count & Judgment \\
\midrule
graph\_color & +8.33pp & +0.00pp & 4/10 & diagnostic only \\
tsp\_mask & +25.00pp & +2.08pp & 2/4 & appendix only \\
countdown & -4.17pp & -4.17pp & 0/2 & negative \\
nqueens & -4.17pp & -10.42pp & 0/2 & negative \\
sat3 & +0.00pp & +0.00pp & 0/4 & near-zero \\
\bottomrule
\end{tabular}
\end{table}


\paragraph{Closure window.}
The final narrow search focused on \texttt{mbpp\_longctx} and \texttt{arc\_mc} before a breadth pivot. The best closure-window results were \texttt{mbpp\_longctx\_seed38\_depthmix\_anchor} at $+5.01$pp and \texttt{arc\_mc\_seed260\_depthmix\_qfirst} at $+9.38$pp, but neither family held up as a clean held-out confirmation line. The strict confirmation queue failed on \texttt{mbpp\_longctx} and flipped negative on \texttt{arc\_mc}.

\paragraph{Breadth pivot.}
Rounds 805--808 were used to stop over-spending on the same \texttt{mbpp\_longctx} recipe family and test a wider set of real tasks. Only \texttt{hotpot} converted to medium-stage positives in that pivot, with a best gain of $+1.00$pp. \texttt{protein\_ss} and \texttt{squad} stayed near the promote gate but did not convert in those breadth rounds.

\section{Reproducibility and Artifact Layout}
The supplementary material is intended to include:
\begin{itemize}
    \item the code for the synthetic task suite under \texttt{rwkv-diff-future-seed/};
    \item the post-training launchers under \texttt{posttrain\_rwkv7/scripts/};
    \item the final queue files used in the closure and breadth pivot;
    \item the synced summaries and JSONL records for rounds 783--808 under \texttt{posttrain\_rwkv7/runs/};
    \item this paper package under \texttt{paper/neurips2025/}.
\end{itemize}
A compact build path is provided in the paper package via \texttt{build.sh}. The tables used in the paper are regenerated from committed paper-side metrics using \texttt{render\_tables.py}.

\section{Assets, Licenses, and Scope}
The real-task audit uses publicly available datasets and benchmark snapshots, including ARC, HotpotQA, SQuAD, MBPP, and a public protein secondary-structure dataset snapshot. We cite the original benchmark creators in the main paper. We do not claim to have completed a full legal audit of every upstream redistribution license inside this submission; this is one reason the checklist answer on explicit license coverage is conservative.


\clearpage
\section*{NeurIPS Paper Checklist}

The checklist is designed to encourage best practices for responsible machine learning research, addressing issues of reproducibility, transparency, research ethics, and societal impact.

\begin{enumerate}

\item {\bf Claims}
    \item[] Question: Do the main claims made in the abstract and introduction accurately reflect the paper's contributions and scope?
    \item[] Answer: \answerYes{}
    \item[] Justification: The abstract and introduction explicitly separate strong synthetic evidence from mixed real-task evidence and do not claim stable held-out confirmation; see Abstract, Sections~1, 4, 5, and 6.

\item {\bf Limitations}
    \item[] Question: Does the paper discuss the limitations of the work performed by the authors?
    \item[] Answer: \answerYes{}
    \item[] Justification: Section~6 discusses the single-GPU budget, the narrow recipe family, the lack of stable held-out real-task confirmation, and the non-i.i.d. nature of the search archive.

\item {\bf Theory assumptions and proofs}
    \item[] Question: For each theoretical result, does the paper provide the full set of assumptions and a complete (and correct) proof?
    \item[] Answer: \answerNA{}
    \item[] Justification: The paper is empirical and does not present formal theorems or proofs.

\item {\bf Experimental result reproducibility}
    \item[] Question: Does the paper fully disclose all the information needed to reproduce the main experimental results of the paper to the extent that it affects the main claims and/or conclusions of the paper (regardless of whether the code and data are provided or not)?
    \item[] Answer: \answerYes{}
    \item[] Justification: Sections~3--5 and Appendix~A--D describe the model family, tasks, queue policy, compute budget, and artifact layout needed to reproduce the reported evidence.

\item {\bf Open access to data and code}
    \item[] Question: Does the paper provide open access to the data and code, with sufficient instructions to faithfully reproduce the main experimental results, as described in supplemental material?
    \item[] Answer: \answerYes{}
    \item[] Justification: The repository snapshot and anonymous supplementary archive include the synthetic task code, post-training launchers, final queue files, synced round-783--808 artifacts, and the paper build scripts; Appendix~D describes the layout.

\item {\bf Experimental setting/details}
    \item[] Question: Does the paper specify all the training and test details (e.g., data splits, hyperparameters, how they were chosen, type of optimizer, etc.) necessary to understand the results?
    \item[] Answer: \answerYes{}
    \item[] Justification: The core paper states the model family, checkpoint, task families, queue gates, and budget ranges, while Appendix~A and D describe the screening policy, artifact layout, and reproduction commands.

\item {\bf Experiment statistical significance}
    \item[] Question: Does the paper report error bars suitably and correctly defined or other appropriate information about the statistical significance of the experiments?
    \item[] Answer: \answerNo{}
    \item[] Justification: We report aggregate medians, positive-count summaries, and some multi-seed synthetic results, but we do not provide a unified error-bar analysis for every real-task family.

\item {\bf Experiments compute resources}
    \item[] Question: For each experiment, does the paper provide sufficient information on the computer resources (type of compute workers, memory, time of execution) needed to reproduce the experiments?
    \item[] Answer: \answerYes{}
    \item[] Justification: Sections~3 and 6 state that the campaign uses a single RTX 4090 and fixed quick-to-medium budgets; Appendix~A clarifies the serial screening policy, common budget ranges, and artifact scope.

\item {\bf Code of ethics}
    \item[] Question: Does the research conducted in the paper conform, in every respect, with the NeurIPS Code of Ethics https://neurips.cc/public/EthicsGuidelines?
    \item[] Answer: \answerYes{}
    \item[] Justification: The work uses public benchmarks, does not involve human subjects, and presents limitations and release considerations conservatively.

\item {\bf Broader impacts}
    \item[] Question: Does the paper discuss both potential positive societal impacts and negative societal impacts of the work performed?
    \item[] Answer: \answerYes{}
    \item[] Justification: Section~6 discusses positive impacts for controllable structured editing and negative risks from misuse of stronger infill or consistency-repair systems.

\item {\bf Safeguards}
    \item[] Question: Does the paper describe safeguards that have been put in place for responsible release of data or models that have a high risk for misuse (e.g., pretrained language models, image generators, or scraped datasets)?
    \item[] Answer: \answerNA{}
    \item[] Justification: This submission does not release a new frontier model or high-risk scraped dataset; it studies a lightweight mechanism on existing model checkpoints and public benchmarks.

\item {\bf Licenses for existing assets}
    \item[] Question: Are the creators or original owners of assets (e.g., code, data, models), used in the paper, properly credited and are the license and terms of use explicitly mentioned and properly respected?
    \item[] Answer: \answerNo{}
    \item[] Justification: We cite the original benchmark creators and describe the assets used, but we do not provide a complete in-paper license audit for every upstream dataset snapshot and dependency.

\item {\bf New assets}
    \item[] Question: Are new assets introduced in the paper well documented and is the documentation provided alongside the assets?
    \item[] Answer: \answerYes{}
    \item[] Justification: The paper and supplement document the new synthetic task suite, queue files, run artifacts, and paper-side generation scripts; see Appendix~D.

\item {\bf Crowdsourcing and research with human subjects}
    \item[] Question: For crowdsourcing experiments and research with human subjects, does the paper include the full text of instructions given to participants and screenshots, if applicable, as well as details about compensation (if any)?
    \item[] Answer: \answerNA{}
    \item[] Justification: The work does not involve crowdsourcing or human-subject data collection.

\item {\bf Institutional review board (IRB) approvals or equivalent for research with human subjects}
    \item[] Question: Does the paper describe potential risks incurred by study participants, whether such risks were disclosed to the subjects, and whether Institutional Review Board (IRB) approvals (or an equivalent approval/review based on the requirements of your country or institution) were obtained?
    \item[] Answer: \answerNA{}
    \item[] Justification: The work does not involve human subjects.

\item {\bf Declaration of LLM usage}
    \item[] Question: Does the paper describe the usage of LLMs if it is an important, original, or non-standard component of the core methods in this research? Note that if the LLM is used only for writing, editing, or formatting purposes and does not impact the core methodology, scientific rigorousness, or originality of the research, declaration is not required.
    \item[] Answer: \answerNA{}
    \item[] Justification: Large language models are not part of the core experimental method studied in this paper.

\end{enumerate}


\end{document}
